\subsection{Comparative Analysis: Empirical Physical Hardware vs. Discrete-Event Simulation}
\label{subsec:sim_vs_real_discussion}

A critical requirement for validating network security frameworks is establishing \textbf{sim-to-real fidelity}—demonstrating that analytical models and discrete-event simulations (e.g., NS-3) accurately predict physical hardware behavior under real-world electromagnetic conditions. Table~\ref{tab_sim_real_comparison_detailed} provides a quantitative comparison between NS-3 simulation predictions and empirical ESP32 hardware testbed measurements across $1,120,555$ telemetry packets.

\begin{table}[htbp]
\caption{Detailed Comparison between NS-3 Simulation and Physical ESP32 Testbed}
\label{tab_sim_real_comparison_detailed}
\centering
\small
\renewcommand{\arraystretch}{1.2}
\resizebox{\columnwidth}{!}{%
\begin{tabular}{|l|c|c|c|}
\hline
\textbf{Performance Metric / Parameter} & \textbf{NS-3 Simulation} & \textbf{Physical ESP32 Testbed} & \textbf{Absolute Delta ($\Delta$)} \\
\hline
Detection Recall / Rate & $98.24\%$ & $98.18\%$ & $\Delta = -0.06\%$ ($\rho = 0.998$) \\
False Positive Rate (FPR) & $0.01\%$ & $0.03\%$ & $\Delta = +0.02\%$ \\
Mean Isolation Latency & $0.85\,\text{s}$ & $0.85\,\text{s}$ & $\Delta = 0.00\,\text{s}$ \\
Normal RSSI Variance & $0.45\,\text{dB}^2$ & $1.21\,\text{dB}^2$ & $\Delta = +0.76\,\text{dB}^2$ \\
Normal IAT Variance & $1.02\,\text{ms}^2$ & $7.84\,\text{ms}^2$ & $\Delta = +6.82\,\text{ms}^2$ \\
Average Energy per Packet & $0.541\,\text{mJ}$ & $0.540\,\text{mJ}$ & $\Delta = -0.001\,\text{mJ}$ \\
Energy Reduction Ratio & $99.40\%$ & $99.40\%$ & Exact Match \\
\hline
\end{tabular}%
}
\end{table}

\subsubsection{Convergent Trends and Sim-to-Real Correlation}
The empirical hardware results strongly validate the theoretical foundation of the \ac{HEDSD} framework, exhibiting near-perfect alignment ($\rho = 0.998$) with simulation trends across three core physical dimensions:
\begin{enumerate}
    \item \textbf{Spatial RSSI Bimodality:} Both simulation and physical hardware experiments demonstrate that physical spatial displacement between Target Node $N_1$ ($3.0\,\text{m}$) and Attacker Node $A_1$ ($8.5\,\text{m}$) induces a sharp bimodal \ac{RSSI} distribution ($-48.2\,\text{dBm}$ vs. $-68.1\,\text{dBm}$), inflating rolling spatial variance beyond the $\sigma^2_{\text{RSSI}} > 18.5\,\text{dB}^2$ threshold.
    \item \textbf{High-Rate Injection Collapse:} Under high-rate DoS bursts (Mode 2), both domains capture a rapid collapse in mean \ac{IAT} ($\bar{\text{IAT}} \downarrow 15.1\,\text{ms}$ in hardware vs. $15.0\,\text{ms}$ in NS-3), triggering instantaneous rate-anomaly flags before sliding window statistics complete.
    \item \textbf{Selective Cryptographic Energy Savings:} Physical INA219 current shunt measurements ($0.540\,\text{mJ/pkt}$) confirmed NS-3 energy predictions ($0.541\,\text{mJ/pkt}$), proving that selective cryptographic auditing eliminates $99.40\%$ of signature verification overhead on resource-constrained microcontrollers.
\end{enumerate}

\subsubsection{Real-World Discrepancies and Noise Factors (Why Physical Results are Less Perfect)}
While the overall detection trends converge, empirical hardware results exhibit minor performance degradation due to physical channel noise not fully modeled in idealized simulation environments:
\begin{enumerate}
    \item \textbf{Ambient Indoor Shadowing and Body Fading ($+0.02\%$ FPR Increase):} In physical indoor corridor environments, transient obstacle shadowing (e.g., personnel walking through the $2.4\,\text{GHz}$ RF propagation path) induces brief $4-6\,\text{dB}$ attenuation dips. Coinciding with minor Wi-Fi transmission queue jitter, these physical fading dips generated $324$ false positive alerts ($0.03\%$ FPR) across $1.08$ million legitimate frames, whereas simulation environments exhibited near-zero false alarms ($0.01\%$ FPR).
    \item \textbf{Microcontroller Crystal Oscillator Thermal Drift ($\Delta \sigma^2_{\text{IAT}} = +6.82\,\text{ms}^2$):} Physical ESP32 microcontrollers suffer from ambient temperature fluctuations ($\pm 40^\circ\text{C}$ thermal variance), causing crystal oscillator drift ($\pm 3.2\,\text{ms}$) and FreeRTOS task switching jitter ($15-50\,\mu\text{s}$). This elevated baseline hardware \ac{IAT} variance from $1.02\,\text{ms}^2$ (simulation) to $7.84\,\text{ms}^2$ (hardware).
    \item \textbf{Receiver AGC Non-Linearity and Quantization Discretization:} Physical ESP32 802.11 Automatic Gain Control (AGC) circuits possess non-linear quantization boundaries near sensitivity limits ($-85\,\text{dBm}$ to $-95\,\text{dBm}$), introducing discrete $1\,\text{dB}$ quantization steps that slightly increase physical RSSI variance ($\Delta \sigma^2_{\text{RSSI}} = +0.76\,\text{dB}^2$).
\end{enumerate}

\subsubsection{Co-located Attacker Failure Mode ($738$ Missed Frames)}
In both simulation and physical hardware, when Attacker $A_1$ is positioned in close physical proximity to Target $N_1$ ($<3.0\,\text{m}$ spatial offset), the spatial RSSI variance drops below the isolation threshold ($\sigma^2_{\text{RSSI}} < 2.0\,\text{dB}^2$), resulting in $738$ missed frames ($1.82\%$ FNR). However, selective cryptographic fallback acts as a secondary defense layer, ensuring that un-isolated frames triggering threat scores above $0.65$ undergo mandatory ECDSA verification, preventing unauthorized network compromise.
